% !TeX root = main.tex
\begin{figure*}[t]
  \centering
  \resizebox{\linewidth}{!}{% Full-width method/pipeline placeholder. Replace with a PDF figure when ready.
\begin{tikzpicture}[
  font=\footnotesize,
  module/.style={
    draw=black!55,
    rounded corners=2pt,
    minimum height=0.58in,
    text width=0.16\textwidth,
    align=center,
    inner sep=5pt
  },
  flow/.style={-{Latex[length=2mm]}, line width=0.75pt, draw=black!70},
  feedback/.style={-{Latex[length=2mm]}, dashed, line width=0.7pt, draw=black!60}
]
  \node[module, fill=blue!7] (input) {\textbf{Inputs}\\Observations, state, or demonstrations};
  \node[module, fill=orange!8, right=0.32in of input] (representation) {\textbf{Representation}\\Encoder or state estimator};
  \node[module, fill=orange!8, right=0.32in of representation] (method) {\textbf{Proposed module}\\Core technical contribution};
  \node[module, fill=orange!8, right=0.32in of method] (decision) {\textbf{Decision}\\Policy, planner, or controller};
  \node[module, fill=green!8, right=0.32in of decision] (output) {\textbf{Outputs}\\Actions and evaluation signals};

  \draw[flow] (input) -- (representation);
  \draw[flow] (representation) -- (method);
  \draw[flow] (method) -- (decision);
  \draw[flow] (decision) -- (output);
  \draw[feedback] (output.south) -- ++(0,-0.22) -| (method.south);
\end{tikzpicture}
}
  \caption{Placeholder for the full method or system pipeline. Replace each
    block with a meaningful stage, show the direction of data flow, and explain
    train-time versus deployment-time components in the caption.}
  \label{fig:pipeline}
\end{figure*}

\subsection{Problem Formulation}

Define the task, assumptions, and notation before describing the method. For a
state $\mathbf{x}_t$ and control $\mathbf{u}_t$, a generic discrete-time model is
\begin{equation}
  \mathbf{x}_{t+1} = f(\mathbf{x}_t, \mathbf{u}_t) + \boldsymbol{\varepsilon}_t,
  \label{eq:dynamics}
\end{equation}
where $\boldsymbol{\varepsilon}_t$ denotes process uncertainty. Refer to
equations by label, as in~\eqref{eq:dynamics}.

\subsection{Approach}

Describe enough implementation detail for another researcher to reproduce
\methodname. Use \figref{fig:pipeline} to give the high-level flow before
introducing individual components.

\subsection{Core Components}

Give each novel component its own paragraph or subsection. For every component,
state its inputs, outputs, training signal, and role in the complete system.
Move secondary derivations to concise prose when space is limited.

\begin{algorithm}[t]
  \caption{Placeholder for \methodname training or inference}
  \label{alg:method}
  \begin{algorithmic}[1]
    \REQUIRE Inputs, model components, and stopping condition
    \STATE Initialize the learnable components and data buffers
    \FOR{each training iteration or control step}
      \STATE Observe the current input
      \STATE Compute the proposed intermediate representation
      \STATE Apply the core method and obtain an output
      \STATE Update the model, controller, or stored data if required
    \ENDFOR
    \RETURN Trained model, policy, plan, or action
  \end{algorithmic}
\end{algorithm}

